%% kapa: exam=1 round=29 year=2018 period=1 form=A subject=civil_law subject_ko=민법 seq=31 number=31
%% source: Q-Net 감정평가사 기출문제 (cst00309) artlSeq=5202808 file=extracted/1008/5202808/2018년 제29회 감정평가사 1차시험 기출문제/2018년 제29회 감정평가사 1차시험 기출문제/1교시 A형.pdf
%% conversion: kordoc 4.24.51 -> tools/gen1.py (auto); review=proofread-sonnet; flags=none
\begin{question}{31}
부동산등기에 관한 설명으로 옳지 않은 것은? (다툼이 있으면 판례에 따름)
\begin{choices}
  \item 멸실된 건물의 보존등기를 신축한 건물의 보존등기로 유용하는 것은 허용되지 않는다.
  \item 소유자로부터 토지를 적법하게 매수한 매수인의 소유권이전등기가 위조된 서류에 의하여 경료되었더라도 그 등기는 유효하다.
  \item 가등기된 권리의 이전등기는 가등기에 대한 부기등기의 형식으로는 경료할 수 없다.
  \item 명의자를 달리하는 중복보존등기가 부동산을 표상함에 부족함이 없는 경우, 선행등기가 원인무효가 아닌 한 후행등기는 실체적 권리관계에 부합하더라도 무효이다.
  \item 지분이전등기가 경료된 경우 그 등기는 적법하게 된 것으로서 진실한 권리상태를 공시하는 것이라고 추정된다.
\end{choices}
\end{question}
\answer{3}
